\documentclass[UTF8, 12pt, fontset=fandol, oneside]{ctexart}
\usepackage{researchnotes}

\graphicspath{{tex/强化学习/images/}}

\title{强化学习：导论}
\author{}
\date{}

\begin{document}

\maketitle

\section*{第 3 章\quad 有限马尔可夫决策过程}

在本章中，我们介绍有限马尔可夫决策过程，简称有限 MDP，的形式化问题；在本书其余部分，我们都试图求解这类问题。
这个问题包含评价性反馈，就像赌博机问题中那样，但它还包含一个关联性方面，即在不同情境中选择不同动作。
MDP 是顺序决策的一种经典形式化方式，在这种决策中，动作不仅影响即时奖励，还会影响后续情境，或者说状态，并通过这些状态影响未来奖励。
因此，MDP 涉及延迟奖励，也涉及在即时奖励与延迟奖励之间进行权衡的需要。
在赌博机问题中，我们估计每个动作 $a$ 的价值 $q_*(a)$；而在 MDP 中，我们估计每个状态 $s$ 中每个动作 $a$ 的价值 $q_*(s,a)$，或者在给定最优动作选择时估计每个状态的价值 $v_*(s)$。
这些依赖于状态的量，对于把长期后果的功劳准确分配给各个动作选择来说是必不可少的。

MDP 是强化学习问题的一种数学理想化形式，适合用来作出精确的理论陈述。
我们将介绍这个问题数学结构中的关键元素，例如回报、价值函数和贝尔曼方程。
我们试图传达的是，可以被表述为有限 MDP 的应用范围非常广。
正如整个人工智能中常见的那样，适用范围的广度与数学上的可处理性之间存在张力。
在本章中，我们介绍这种张力，并讨论它所带来的一些取舍和挑战。
强化学习超越 MDP 的若干方式将在第 17 章讨论。

\subsection*{3.1\quad 智能体--环境接口}

MDP 旨在把通过交互学习以达成目标的问题，框定为一种直接明了的形式。
学习者和决策者称为智能体。
与智能体交互的对象，也就是智能体之外的一切，称为环境。
二者持续交互：智能体选择动作，环境对这些动作作出回应，并向智能体呈现新的情境。
环境还会产生奖励；奖励是特殊的数值，智能体希望通过自己的动作选择，在时间上最大化这些数值。

\begin{figure}[htbp]
  \centering
  \includegraphics[width=0.74\textwidth]{rl4-fig-3-1.png}
  \caption{马尔可夫决策过程中的智能体--环境交互。}
\end{figure}

更具体地说，智能体与环境在一串离散时间步 $t=0,1,2,3,\ldots$ 上逐步交互。
在每个时间步 $t$，智能体接收环境状态的某种表示 $S_t\in\mathcal{S}$，并在此基础上选择一个动作 $A_t\in\mathcal{A}(s)$。
一个时间步之后，部分地作为其动作的结果，智能体接收到一个数值奖励 $R_{t+1}\in\mathcal{R}\subset\mathbb{R}$，并发现自己处在一个新的状态 $S_{t+1}$ 中。
于是，MDP 与智能体共同产生一个序列，或称轨迹，其开头如下：
\begin{equation}
S_0,A_0,R_1,S_1,A_1,R_2,S_2,A_2,R_3,\ldots
\tag{3.1}
\end{equation}

在有限 MDP 中，状态、动作和奖励的集合（$\mathcal{S}$、$\mathcal{A}$ 和 $\mathcal{R}$）都含有有限个元素。
在这种情况下，随机变量 $R_t$ 和 $S_t$ 具有定义良好的离散概率分布，并且该分布只依赖于前一个状态和前一个动作。
也就是说，对于这些随机变量的特定取值 $s'\in\mathcal{S}$ 和 $r\in\mathcal{R}$，在给定前一个状态和动作的特定取值时，它们在时刻 $t$ 出现的概率为
\begin{equation}
p(s',r|s,a)\doteq \Pr\{S_t=s',R_t=r\mid S_{t-1}=s,A_{t-1}=a\},
\tag{3.2}
\end{equation}
其中对所有 $s',s\in\mathcal{S}$、$r\in\mathcal{R}$ 和 $a\in\mathcal{A}(s)$ 都成立。
函数 $p$ 定义了 MDP 的动态。
等号上方的点提醒我们，这是一个定义（这里是函数 $p$ 的定义），而不是从先前定义推出的事实。
动态函数 $p:\mathcal{S}\times\mathcal{R}\times\mathcal{S}\times\mathcal{A}\to[0,1]$ 是一个普通的、具有四个自变量的确定性函数。
其中间的竖线“$|$”来自条件概率的记号；但在这里，它只是提醒我们，对于每一组 $s$ 和 $a$ 的选择，$p$ 指定了一个概率分布，也就是说
\begin{equation}
\sum_{s'\in\mathcal{S}}\sum_{r\in\mathcal{R}}p(s',r|s,a)=1,\quad \text{对所有 }s\in\mathcal{S},\ a\in\mathcal{A}(s).
\tag{3.3}
\end{equation}

在马尔可夫决策过程中，由 $p$ 给出的概率完全刻画了环境的动态。
也就是说，$S_t$ 和 $R_t$ 的每个可能取值的概率，只依赖于紧接在它们之前的状态和动作 $S_{t-1}$ 与 $A_{t-1}$；在给定它们之后，完全不依赖更早的状态和动作。
最好把这一点看作对状态的限制，而不是对决策过程本身的限制。
状态必须包含过去智能体--环境交互中所有会对未来产生影响的信息。
如果状态确实包含这些信息，那么就说该状态具有马尔可夫性质。
从本书这里开始，我们将假定马尔可夫性质成立；不过在第二部分中，我们会考虑不依赖它的近似方法，而在第 17 章中，我们会考虑如何从非马尔可夫观测中学习和构造马尔可夫状态。

由四自变量动态函数 $p$ 出发，可以计算出我们可能希望知道的关于环境的其他任何量，例如状态转移概率。
我们稍微滥用记号，把状态转移概率记为三自变量函数 $p:\mathcal{S}\times\mathcal{S}\times\mathcal{A}\to[0,1]$：
\begin{equation}
p(s'|s,a)\doteq \Pr\{S_t=s'\mid S_{t-1}=s,A_{t-1}=a\}
=\sum_{r\in\mathcal{R}}p(s',r|s,a).
\tag{3.4}
\end{equation}
我们还可以把状态--动作对的期望奖励计算为二自变量函数 $r:\mathcal{S}\times\mathcal{A}\to\mathbb{R}$：
\begin{equation}
r(s,a)\doteq \mathbb{E}[R_t\mid S_{t-1}=s,A_{t-1}=a]
=\sum_{r\in\mathcal{R}}r\sum_{s'\in\mathcal{S}}p(s',r|s,a),
\tag{3.5}
\end{equation}
并且可以把状态--动作--下一状态三元组的期望奖励计算为三自变量函数 $r:\mathcal{S}\times\mathcal{A}\times\mathcal{S}\to\mathbb{R}$：
\begin{equation}
r(s,a,s')\doteq \mathbb{E}[R_t\mid S_{t-1}=s,A_{t-1}=a,S_t=s']
=\sum_{r\in\mathcal{R}}r\frac{p(s',r|s,a)}{p(s'|s,a)}.
\tag{3.6}
\end{equation}
在本书中，我们通常使用式（3.2）中的四自变量 $p$ 函数，但这些其他记号有时也很方便。

MDP 框架是抽象而灵活的，可以以许多不同方式应用到许多不同问题。
例如，时间步不必指固定的真实时间间隔；它们可以指决策和行动的任意连续阶段。
动作可以是低层控制，例如施加到机器人手臂电机上的电压，也可以是高层决策，例如是否吃午饭，或者是否去读研究生。
类似地，状态也可以有多种形式。
状态可以完全由低层感觉决定，例如直接的传感器读数；也可以更高层、更抽象，例如对房间中物体的符号描述。
组成状态的一部分内容可以基于过去感觉的记忆，甚至可以完全是心理的或主观的。
例如，一个智能体可能处在“不确定某个物体在哪里”的状态，或者处在“刚刚在某种明确定义的意义上受到惊讶”的状态。
同样，有些动作可能完全是心理性的或计算性的。
例如，有些动作可能控制智能体选择思考什么，或者把注意力集中在哪里。
一般来说，动作可以是我们希望学习如何作出的任何决策，而状态可以是任何我们知道、且可能有助于作出这些决策的信息。

尤其要注意的是，智能体与环境的边界通常并不等同于机器人或动物身体的物理边界。
通常，这条边界会画得比身体边界更靠近智能体。
例如，机器人的电机、机械连杆以及感知硬件通常应被视为环境的一部分，而不是智能体的一部分。
类似地，如果我们把 MDP 框架应用于人或动物，那么肌肉、骨骼和感觉器官应被视为环境的一部分。
奖励也一样，它们大概是在自然或人工学习系统的物理身体内部计算出来的，但仍被视为智能体外部的东西。

我们遵循的一般规则是：凡是智能体不能任意改变的东西，都被认为位于智能体之外，因而是其环境的一部分。
我们并不假设环境中的一切对智能体来说都是未知的。
例如，智能体常常相当了解其奖励如何作为动作以及采取动作时所在状态的函数而被计算出来。
但是，我们总是把奖励计算视为智能体外部的东西，因为它定义了智能体面对的任务，因此必须超出智能体任意改变的能力。
事实上，在某些情况下，智能体可能知道关于环境如何运作的一切，却仍然面对一个困难的强化学习任务；这就像我们也许准确知道魔方如何运作，却仍然无法把它复原一样。
智能体--环境边界表示的是智能体绝对控制的限度，而不是它知识的限度。

智能体--环境边界可以为了不同目的而放在不同位置。
在一个复杂机器人中，可能同时运行着许多不同智能体，每个智能体都有自己的边界。
例如，一个智能体可以作出高层决策，而这些决策构成另一个较低层智能体所面对状态的一部分，后者负责实现这些高层决策。
在实践中，一旦我们选择了特定的状态、动作和奖励，从而确定了一个感兴趣的具体决策任务，智能体--环境边界也就随之确定。

MDP 框架是对通过交互进行目标导向学习的问题的一种相当大的抽象。
它提出，无论感觉、记忆和控制装置的细节如何，无论要实现的目标是什么，任何学习目标导向行为的问题都可以归约为智能体与环境之间来回传递的三个信号：一个信号表示智能体作出的选择（动作），一个信号表示这些选择所依据的基础（状态），一个信号定义智能体的目标（奖励）。
这个框架未必足以有效表示所有决策学习问题，但已经证明它非常有用且适用广泛。

当然，具体的状态和动作会随着任务而有很大差异，而它们的表示方式会强烈影响性能。
在强化学习中，正如在其他类型的学习中一样，目前这类表示选择更多是一门艺术，而不是科学。
在本书中，我们会给出一些关于状态和动作良好表示方式的建议和例子，但我们的主要关注点，是在表示已经选定之后，用于学习如何行动的一般原则。

\noindent\textbf{例 3.1：生物反应器}\quad 假设强化学习被用于逐时刻地确定生物反应器的温度和搅拌速率；生物反应器是一个装有营养物和细菌的大容器，用来生产有用的化学品。
在这样的应用中，动作可能是目标温度和目标搅拌速率，它们被传递给低层控制系统，而低层控制系统又直接启动加热元件和电机以达到这些目标。
状态很可能是热电偶和其他传感器读数，也许经过滤波和延迟，再加上表示容器中成分和目标化学品的符号输入。
奖励可能是生物反应器生产有用化学品的速率在每一时刻的度量。
注意，在这里每个状态都是由传感器读数和符号输入组成的列表，或向量；每个动作也是由目标温度和搅拌速率组成的向量。
强化学习任务中的状态和动作具有这种结构化表示是很典型的。
另一方面，奖励始终是单个数。

\noindent\textbf{例 3.2：抓取--放置机器人}\quad 考虑用强化学习来控制机器人手臂在一个重复的抓取--放置任务中的运动。
如果我们想学习快速而平滑的运动，那么学习智能体将必须直接控制电机，并且获得关于机械连杆当前位置和速度的低延迟信息。
在这种情况下，动作可能是在每个关节处施加到各个电机的电压，而状态可能是关节角和速度的最新读数。
奖励可以是每成功抓起并放置一个物体给 $+1$。
为了鼓励平滑运动，在每个时间步还可以根据运动逐时刻的“不平滑程度”给出一个小的负奖励。

\noindent\textbf{练习 3.1}\quad 自己设计三个适合 MDP 框架的示例任务，并为每个任务指出它的状态、动作和奖励。
让这三个例子彼此尽可能不同。
这个框架抽象而灵活，可以以许多不同方式应用。
在你的至少一个例子中，试着以某种方式拓展它的边界。

\noindent\textbf{练习 3.2}\quad MDP 框架是否足以有效表示所有目标导向的学习任务？
你能想到任何明确的例外吗？

\noindent\textbf{练习 3.3}\quad 考虑驾驶问题。
你可以用油门、方向盘和刹车来定义动作，也就是在你的身体与机器接触的地方定义动作。
或者你可以把动作定义得更靠外一些，例如在橡胶与道路接触的地方，把你的动作看作轮胎扭矩。
又或者，你可以把动作定义得更靠内一些，例如在你的大脑与身体接触的地方，把动作看作控制肢体的肌肉抽动。
你甚至可以上升到很高的层次，说你的动作是你选择开车去哪里。
正确的层次是什么，智能体与环境之间的边界应该画在哪里？
在什么依据上，一个边界位置应优先于另一个？
是否存在某种根本性理由来偏好某一个位置，还是说这只是一个自由选择？

\noindent\textbf{例 3.3：回收机器人}\quad 一个移动机器人负责在办公室环境中收集空汽水罐。
它有用于检测汽水罐的传感器，还有能够把汽水罐拾起并放入自身储物箱的机械臂和夹爪；它依靠可充电电池运行。
机器人的控制系统包含若干组件，用于解释传感信息、导航，以及控制机械臂和夹爪。
关于如何搜索汽水罐的高层决策，由一个强化学习智能体根据电池当前电量作出。
为了构造一个简单例子，我们假定只区分两个电量水平，并由此得到小的状态集 $\mathcal{S}=\{\text{high},\text{low}\}$。
在每个状态中，智能体可以决定：（1）主动搜索一段时间以寻找汽水罐；（2）保持不动，等待某人给它带来汽水罐；或者（3）返回基地为电池充电。
当电量为 \texttt{high} 时，充电总是愚蠢的，因此我们不把它包含在该状态的动作集中。
于是动作集为
\[
\mathcal{A}(\text{high})=\{\text{search},\text{wait}\},\qquad
\mathcal{A}(\text{low})=\{\text{search},\text{wait},\text{recharge}\}.
\]

奖励在大多数时候为零，但当机器人获得一个空罐时为正，而当电池完全耗尽时为一个很大的负值。
寻找汽水罐的最好方式是主动搜索，但这会耗费机器人的电量，而等待则不会。
只要机器人正在搜索，它的电池就有可能被耗尽。
在这种情况下，机器人必须关机并等待救援，从而产生一个低奖励。
如果电量为 \texttt{high}，那么一次主动搜索总能完成，而没有耗尽电池的风险。
一次从 \texttt{high} 电量开始的搜索，会以概率 $\alpha$ 使电量仍为 \texttt{high}，并以概率 $1-\alpha$ 使电量降为 \texttt{low}。
另一方面，当电量为 \texttt{low} 时进行搜索，会以概率 $\beta$ 使电量保持 \texttt{low}，并以概率 $1-\beta$ 耗尽电池。
在后一种情况下，机器人必须被救援，然后电池被重新充至 \texttt{high}。
机器人收集到的每个汽水罐都算作一个单位奖励，而每当机器人必须被救援时，奖励为 $-3$。
令 $r_{\text{search}}$ 和 $r_{\text{wait}}$ 分别表示机器人在搜索和等待时将收集到的汽水罐的期望数量（因而也是期望奖励），并且 $r_{\text{search}}>r_{\text{wait}}$。
最后，假设在返回基地充电途中无法收集汽水罐，并且在电池耗尽的那个时间步也无法收集汽水罐。
这个系统于是是一个有限 MDP；我们可以写出它的转移概率和期望奖励，其动态如下图左侧表格所示：

\begin{figure}[htbp]
  \centering
  \includegraphics[width=0.86\textwidth]{rl4-recycling-robot-dynamics.png}
  \caption{回收机器人示例中的转移表与转移图。}
\end{figure}

注意，表中每一行对应当前状态 $s$、动作 $a\in\mathcal{A}(s)$ 和下一状态 $s'$ 的一种可能组合。
有些转移发生的概率为零，因此没有为它们指定期望奖励。
右侧所示是另一种概括有限 MDP 动态的有用方式，称为转移图。
图中有两类节点：状态节点和动作节点。
每个可能状态对应一个状态节点，即用状态名称标记的大空心圆；每个状态--动作对对应一个动作节点，即用动作名称标记的小实心圆，并由一条线连接到相应状态节点。
从状态 $s$ 出发并采取动作 $a$，会沿着从状态节点 $s$ 到动作节点 $(s,a)$ 的线移动。
然后，环境经由离开动作节点 $(s,a)$ 的某条箭头，对下一状态节点作出转移响应。
每条箭头对应一个三元组 $(s,s',a)$，其中 $s'$ 是下一状态；我们用转移概率 $p(s'|s,a)$ 和该转移的期望奖励 $r(s,a,s')$ 来标记箭头。
注意，离开一个动作节点的各箭头上标注的转移概率之和总是 $1$。

\noindent\textbf{练习 3.4}\quad 给出一个与例 3.3 中表格类似的表，但针对 $p(s',r|s,a)$。
它应包含 $s$、$a$、$s'$、$r$ 和 $p(s',r|s,a)$ 各列，并且对每个满足 $p(s',r|s,a)>0$ 的四元组都有一行。

\subsection*{3.2\quad 目标与奖励}

在强化学习中，智能体的目的或目标通过一个特殊信号来形式化，这个信号称为奖励，并从环境传递给智能体。
在每个时间步，奖励是一个简单的数 $R_t\in\mathbb{R}$。
非正式地说，智能体的目标是最大化它接收到的奖励总量。
这意味着要最大化的不是即时奖励，而是长远来看累积起来的奖励。
我们可以把这个非正式想法清楚地表述为奖励假设：

\begin{quote}
我们所谓的目标和目的，都可以很好地理解为：最大化一个接收到的标量信号（称为奖励）的累积和的期望值。
\end{quote}

使用奖励信号来形式化目标这一想法，是强化学习最鲜明的特征之一。
虽然用奖励信号来表述目标乍看之下可能显得受限，但在实践中，它已经证明自己灵活且适用广泛。
理解这一点的最好方式，是考虑它已经如何被使用，或者可以如何被使用的例子。
例如，为了让机器人学会行走，研究者可以在每个时间步给予一个与机器人向前运动成比例的奖励。
在让机器人学习如何逃出迷宫时，奖励常常是在逃出之前经过的每个时间步为 $-1$；这会鼓励智能体尽可能快地逃出。
为了让机器人学会寻找并收集空汽水罐以便回收，可以在大多数时候给它零奖励，然后对每个收集到的汽水罐给 $+1$ 奖励。
也可以在机器人撞到东西或有人对它大喊时给它负奖励。
对于一个学习下跳棋或国际象棋的智能体，赢棋时的自然奖励是 $+1$，输棋时是 $-1$，和棋以及所有非终止局面都是 $0$。

你可以看出这些例子中发生了什么。
智能体总是在学习最大化自己的奖励。
如果我们希望它为我们做某件事，就必须以这样一种方式为它提供奖励：当智能体最大化这些奖励时，它同时也会实现我们的目标。
因此，我们设置的奖励真正指示我们希望完成什么，这一点至关重要。
尤其是，奖励信号并不是向智能体传授关于如何达成目标的先验知识的地方。
例如，一个下棋智能体应当只因为实际获胜而得到奖励，而不是因为达成诸如吃掉对方棋子或控制棋盘中心这样的子目标而得到奖励。
如果这类子目标得到了奖励，那么智能体可能找到一种达成子目标却没有达成真正目标的方法。
例如，它可能找到一种吃掉对方棋子的方法，却以输掉比赛为代价。
奖励信号是你向机器人传达“你希望它实现什么”的方式，而不是传达“你希望它如何实现”的方式。

\subsection*{3.3\quad 回报与幕}

到目前为止，我们只是非正式地讨论了学习目标。
我们说过，智能体的目标是在长远上最大化它接收到的累积奖励。
这如何形式化定义呢？
如果把时间步 $t$ 之后接收到的奖励序列记为 $R_{t+1},R_{t+2},R_{t+3},\ldots$，那么我们究竟希望最大化这个序列的哪一方面？
一般来说，我们希望最大化期望回报；回报记为 $G_t$，定义为奖励序列的某个特定函数。
在最简单的情况下，回报就是奖励之和：
\begin{equation}
G_t\doteq R_{t+1}+R_{t+2}+R_{t+3}+\cdots+R_T,
\tag{3.7}
\end{equation}
其中 $T$ 是最终时间步。
这种方式适用于存在自然最终时间步概念的应用，也就是智能体--环境交互会自然分解为若干子序列的情形；这些子序列称为幕，例如一局游戏、一次穿越迷宫，或任何类型的重复交互。
每一幕都在一个特殊状态中结束，这个状态称为终止状态；随后系统会重置到某个标准初始状态，或者从一个标准初始状态分布中采样。
即使你把各幕看作以不同方式结束，例如赢得或输掉一局游戏，下一幕的开始也与前一幕如何结束无关。
因此，这些幕都可以被看作在同一个终止状态中结束，只是不同结果对应不同奖励。
具有这种幕的任务称为分幕任务。
在分幕任务中，我们有时需要把所有非终止状态组成的集合 $\mathcal{S}$，与所有状态加上终止状态组成的集合 $\mathcal{S}^+$ 区分开来。
终止时间 $T$ 是一个随机变量，通常会因幕而异。

另一方面，在许多情况下，智能体--环境交互不会自然分解为可识别的幕，而是无界地持续下去。
例如，对于一个持续运行的过程控制任务，或者一个寿命很长的机器人应用，这会是自然的表述方式。
我们把这类任务称为持续任务。
对持续任务来说，回报定义（3.7）是有问题的，因为最终时间步将是 $T=\infty$，而作为我们试图最大化对象的回报本身也很容易是无限的。
例如，假设智能体在每个时间步都收到 $+1$ 的奖励。
因此，本书中我们通常使用一个概念上稍微复杂、但数学上简单得多的回报定义。

我们需要的额外概念是折扣。
按照这种方法，智能体试图选择动作，使它未来收到的折扣奖励之和最大化。
具体地说，它选择 $A_t$ 以最大化期望折扣回报：
\begin{equation}
G_t\doteq R_{t+1}+\gamma R_{t+2}+\gamma^2R_{t+3}+\cdots
=\sum_{k=0}^{\infty}\gamma^kR_{t+k+1},
\tag{3.8}
\end{equation}
其中 $\gamma$ 是一个参数，$0\leq\gamma\leq1$，称为折扣率。
折扣率决定未来奖励的现值：未来 $k$ 个时间步后收到的奖励，只值其立即收到时价值的 $\gamma^{k-1}$ 倍。
如果 $\gamma<1$，只要奖励序列 $\{R_k\}$ 有界，式（3.8）中的无穷和就是有限值。
如果 $\gamma=0$，智能体就只关心最大化即时奖励，因而是“近视的”：此时它的目标是学习如何选择 $A_t$，使得只有 $R_{t+1}$ 被最大化。
如果智能体的每个动作都只影响即时奖励，而不影响未来奖励，那么近视智能体可以通过分别最大化每个即时奖励来最大化（3.8）。
但一般来说，为了最大化即时奖励而行动，会降低对未来奖励的获取，从而降低回报。
当 $\gamma$ 接近 $1$ 时，回报目标会更强地考虑未来奖励，智能体也变得更有远见。

相邻时间步上的回报之间有一种关系，这对强化学习理论和算法都很重要：
\begin{equation}
\begin{aligned}
G_t &\doteq R_{t+1}+\gamma R_{t+2}+\gamma^2R_{t+3}+\gamma^3R_{t+4}+\cdots\\
&=R_{t+1}+\gamma\left(R_{t+2}+\gamma R_{t+3}+\gamma^2R_{t+4}+\cdots\right)\\
&=R_{t+1}+\gamma G_{t+1}.
\end{aligned}
\tag{3.9}
\end{equation}
注意，只要我们定义 $G_T=0$，这个关系对所有 $t<T$ 的时间步都成立，即使终止发生在 $t+1$。
这通常使得从奖励序列计算回报变得容易。

还要注意，虽然回报（3.8）是无穷多项的和，但当奖励非零且恒定时，只要 $\gamma<1$，它仍然是有限的。
例如，如果奖励恒为 $+1$，那么回报为
\begin{equation}
G_t=\sum_{k=0}^{\infty}\gamma^k=\frac{1}{1-\gamma}.
\tag{3.10}
\end{equation}

\noindent\textbf{练习 3.5}\quad 第 3.1 节中的方程是针对持续情形的；若要应用于分幕任务，需要进行很小的修改。
请通过给出（3.3）的修改版本来说明你知道需要怎样修改。

\noindent\textbf{例 3.4：倒立摆}\quad 这个任务的目标是向一辆沿轨道运动的小车施加力，使铰接在小车上的杆保持不倒。

\begin{figure}[htbp]
  \centering
  \includegraphics[width=0.52\textwidth]{rl4-pole-balancing.png}
  \caption{倒立摆任务示意图。}
\end{figure}

如果杆偏离竖直方向超过给定角度，或者小车跑出轨道，就认为发生失败。
每次失败后，杆被重置为竖直。
这个任务可以被看作分幕任务，其中自然的幕就是反复尝试平衡杆。
在这种情况下，奖励可以是在没有发生失败的每个时间步给 $+1$，因此每个时间步的回报就是直到失败为止的步数。
在这种情况下，永远成功保持平衡意味着回报为无穷。
或者，我们也可以把倒立摆看作持续任务，并使用折扣。
在这种情况下，每次失败时奖励为 $-1$，所有其他时间奖励为零。
于是每个时间步的回报与 $-\gamma^K$ 有关，其中 $K$ 是失败前的时间步数。
无论采用哪种情况，回报都是通过尽可能长时间地保持杆的平衡而最大化的。

\noindent\textbf{练习 3.6}\quad 假设你把倒立摆作为分幕任务处理，但也使用折扣，并且除失败时为 $-1$ 外，所有奖励都为零。
那么每个时间步的回报是什么？
这个回报与该任务的折扣持续表述中的回报有何不同？

\noindent\textbf{练习 3.7}\quad 想象你正在设计一个跑迷宫的机器人。
你决定给它逃出迷宫时 $+1$ 的奖励，并在所有其他时间给零奖励。
这个任务似乎会自然地分解为若干幕，也就是一次次穿越迷宫；因此你决定把它作为分幕任务处理，目标是最大化期望总奖励（3.7）。
运行学习智能体一段时间后，你发现它在逃出迷宫方面没有表现出任何改进。
出了什么问题？
你是否有效地向智能体传达了你希望它实现什么？

\noindent\textbf{练习 3.8}\quad 假设 $\gamma=0.5$，并且收到如下奖励序列：$R_1=-1$、$R_2=2$、$R_3=6$、$R_4=3$、$R_5=2$，且 $T=5$。
$G_0,G_1,\ldots,G_5$ 分别是多少？
提示：向后计算。

\noindent\textbf{练习 3.9}\quad 假设 $\gamma=0.9$，奖励序列为 $R_1=2$，随后是无限多个 $7$。
$G_1$ 和 $G_0$ 分别是多少？

\noindent\textbf{练习 3.10}\quad 证明（3.10）中的第二个等号。

\subsection*{3.4\quad 分幕任务与持续任务的统一记号}

在前一节中，我们描述了两类强化学习任务：一类中智能体--环境交互会自然地分解为一系列彼此独立的幕，即分幕任务；另一类中这种分解不会发生，即持续任务。
前一种情况在数学上更容易，因为每个动作只影响随后在该幕中接收到的有限多个奖励。
在本书中，我们有时考虑一类问题，有时考虑另一类问题，但也常常同时考虑二者。
因此，建立一种能够让我们同时精确谈论两种情况的记号是有用的。

为了精确描述分幕任务，还需要一些额外记号。
我们不再考虑一条很长的时间步序列，而需要考虑一系列幕，每一幕都由有限个时间步组成。
每一幕的时间步都从零重新编号。
因此，我们不仅要引用 $S_t$，即时间 $t$ 的状态表示，还要引用 $S_{t,i}$，即第 $i$ 幕中时间 $t$ 的状态表示；对 $A_{t,i}$、$R_{t,i}$、$\pi_{t,i}$、$T_i$ 等也是如此。
不过，事实证明，当我们讨论分幕任务时，几乎从不需要区分不同的幕。
我们几乎总是在考虑某一个特定的单幕，或者陈述对所有幕都成立的事情。
因此，在实践中，我们几乎总是略微滥用记号，省略对幕编号的显式引用。
也就是说，我们写 $S_t$ 来表示 $S_{t,i}$，其他量也类似。

为了得到一个同时涵盖分幕任务和持续任务的单一记号，我们还需要一个约定。
我们已经在一种情况下把回报定义为有限项之和（3.7），在另一种情况下定义为无限项之和（3.8）。
这两者可以通过把一幕的终止看作进入一个特殊的吸收状态而统一起来；该状态只转移到自身，并且只产生零奖励。
例如，考虑下面的状态转移图：

\begin{figure}[htbp]
  \centering
  \includegraphics[width=0.64\textwidth]{rl4-absorbing-state.png}
  \caption{用吸收状态统一分幕任务与持续任务的记号。}
\end{figure}

这里，实心方块表示对应于一幕结束的特殊吸收状态。
从 $S_0$ 开始，我们得到奖励序列 $+1,+1,+1,0,0,0,\ldots$。
对它们求和时，无论只对前 $T$ 个奖励求和（这里 $T=3$），还是对完整无限序列求和，得到的回报都相同。
即使引入折扣，这一点仍然成立。
因此，一般地，我们可以按（3.8）定义回报，在不需要时省略幕编号，并且允许在总和有定义时取 $\gamma=1$，例如因为所有幕都会终止。
或者，我们也可以写为
\begin{equation}
G_t\doteq \sum_{k=t+1}^{T}\gamma^{k-t-1}R_k,
\tag{3.11}
\end{equation}
其中允许 $T=\infty$ 或 $\gamma=1$，但二者不能同时发生。
在本书其余部分，我们都使用这些约定来简化记号，并表达分幕任务与持续任务之间的紧密平行。
稍后在第 10 章中，我们会介绍一种既是持续的、又不使用折扣的表述。

\subsection*{3.5\quad 策略与价值函数}

几乎所有强化学习算法都涉及估计价值函数。
价值函数是关于状态（或状态--动作对）的函数，它们估计智能体处在某个给定状态中有多好，或者在某个给定状态中执行某个给定动作有多好。
这里“多好”的概念是用未来可以期望得到的奖励来定义的；更精确地说，是用期望回报来定义的。
当然，智能体未来能够期望收到的奖励，取决于它将采取什么动作。
因此，价值函数是相对于特定行动方式定义的，这些行动方式称为策略。

形式上，策略是从状态到选择每个可能动作的概率的映射。
如果智能体在时间 $t$ 遵循策略 $\pi$，那么当 $S_t=s$ 时，$\pi(a|s)$ 就是 $A_t=a$ 的概率。
和 $p$ 一样，$\pi$ 是一个普通函数；$\pi(a|s)$ 中间的“$|$”只是提醒我们，对每个 $s\in\mathcal{S}$，它在 $a\in\mathcal{A}(s)$ 上定义了一个概率分布。
强化学习方法规定了智能体的策略如何由于其经验而改变。

\noindent\textbf{练习 3.11}\quad 如果当前状态是 $S_t$，并且动作按随机策略 $\pi$ 选择，那么用 $\pi$ 和四自变量函数 $p$（3.2）如何表示 $R_{t+1}$ 的期望？

策略 $\pi$ 下状态 $s$ 的价值函数记为 $v_\pi(s)$，它是在 $s$ 中开始并随后遵循 $\pi$ 时的期望回报。
对于 MDP，我们可以形式化地定义 $v_\pi$：
\begin{equation}
v_\pi(s)\doteq \mathbb{E}_\pi[G_t\mid S_t=s]
=\mathbb{E}_\pi\!\left[\sum_{k=0}^{\infty}\gamma^kR_{t+k+1}\mid S_t=s\right],
\quad \text{对所有 }s\in\mathcal{S}.
\tag{3.12}
\end{equation}
其中 $\mathbb{E}_\pi[\cdot]$ 表示在智能体遵循策略 $\pi$ 的条件下随机变量的期望，而 $t$ 是任意时间步。
注意，如果存在终止状态，则终止状态的价值总是零。
我们把函数 $v_\pi$ 称为策略 $\pi$ 的状态价值函数。

类似地，我们把在策略 $\pi$ 下，处于状态 $s$ 并采取动作 $a$ 的价值记为 $q_\pi(s,a)$；它是在 $s$ 中开始、采取动作 $a$，并随后遵循策略 $\pi$ 时的期望回报：
\begin{equation}
q_\pi(s,a)\doteq \mathbb{E}_\pi[G_t\mid S_t=s,A_t=a]
=\mathbb{E}_\pi\!\left[\sum_{k=0}^{\infty}\gamma^kR_{t+k+1}\mid S_t=s,A_t=a\right].
\tag{3.13}
\end{equation}
我们把 $q_\pi$ 称为策略 $\pi$ 的动作价值函数。

\noindent\textbf{练习 3.12}\quad 给出一个用 $q_\pi$ 和 $\pi$ 表示 $v_\pi$ 的方程。

\noindent\textbf{练习 3.13}\quad 给出一个用 $v_\pi$ 和四自变量 $p$ 表示 $q_\pi$ 的方程。

价值函数 $v_\pi$ 和 $q_\pi$ 可以从经验中估计。
例如，如果一个智能体遵循策略 $\pi$，并且对每个遇到的状态，都维护该状态之后实际回报的平均值，那么随着该状态被遇到次数趋于无穷，这个平均值将收敛到该状态的价值 $v_\pi(s)$。
如果对每个状态中采取的每个动作分别维护平均值，那么这些平均值同样会收敛到动作价值 $q_\pi(s,a)$。
我们把这类估计方法称为蒙特卡洛方法，因为它们涉及对许多实际回报的随机样本求平均。
这类方法将在第 5 章介绍。
当然，如果状态非常多，那么为每个状态单独维护平均值可能并不实际。
相反，智能体必须把 $v_\pi$ 和 $q_\pi$ 维护为带参数的函数（参数个数少于状态个数），并调节这些参数，使其更好地匹配观测到的回报。
这也可以产生准确估计，尽管很大程度上取决于参数化函数近似器的性质。
这些可能性将在本书第二部分讨论。

贯穿强化学习和动态规划，所使用的价值函数有一个基本性质：它们满足类似于我们已经为回报建立的递归关系（3.9）。
对于任意策略 $\pi$ 和任意状态 $s$，状态 $s$ 的价值与其可能后继状态的价值之间满足如下相容性条件：
\begin{equation}
\begin{aligned}
v_\pi(s)
&\doteq \mathbb{E}_\pi[G_t\mid S_t=s]\\
&= \mathbb{E}_\pi[R_{t+1}+\gamma G_{t+1}\mid S_t=s]\qquad \text{（由（3.9））}\\
&= \sum_a\pi(a|s)\sum_{s'}\sum_r p(s',r|s,a)
\left[r+\gamma\mathbb{E}_\pi[G_{t+1}\mid S_{t+1}=s']\right]\\
&= \sum_a\pi(a|s)\sum_{s',r}p(s',r|s,a)\left[r+\gamma v_\pi(s')\right],
\quad \text{对所有 }s\in\mathcal{S}.
\end{aligned}
\tag{3.14}
\end{equation}
这里隐含的是，动作 $a$ 取自集合 $\mathcal{A}(s)$，下一状态 $s'$ 取自集合 $\mathcal{S}$（在分幕问题中取自 $\mathcal{S}^+$），奖励 $r$ 取自集合 $\mathcal{R}$。
还要注意，最后一个方程把两个求和合并了：一个遍历 $s'$ 的所有取值，另一个遍历 $r$ 的所有取值，合并为对二者所有可能取值的一次求和。
我们经常使用这种合并求和来简化公式。
还要注意，最后的表达式很容易读成一个期望。
它实际上是在对三个变量 $a$、$s'$ 和 $r$ 的所有取值求和。
对每个三元组，我们计算它的概率 $\pi(a|s)p(s',r|s,a)$，用这个概率加权括号中的量，然后对所有可能性求和得到期望。

方程（3.14）是 $v_\pi$ 的贝尔曼方程。
它表达了一个状态的价值与其后继状态价值之间的关系。
可以把它想象为从某个状态向前看其可能的后继状态，如下方备份图所示。

\begin{figure}[htbp]
  \centering
  \includegraphics[width=0.26\textwidth]{rl4-backup-vpi.png}
  \caption{$v_\pi$ 的备份图。}
\end{figure}

每个空心圆表示一个状态，每个实心圆表示一个状态--动作对。
从状态 $s$，也就是顶部的根节点开始，智能体可以根据其策略 $\pi$ 采取若干动作之一；图中画出了三个。
从每个动作出发，环境可能根据函数 $p$ 给出的动态，以若干下一状态 $s'$（图中画出两个）和一个奖励 $r$ 作出响应。
贝尔曼方程（3.14）对所有可能性取平均，并按其发生概率加权。
它说明，起始状态的价值必须等于期望下一状态的折扣价值，加上沿途期望得到的奖励。

价值函数 $v_\pi$ 是其贝尔曼方程的唯一解。
在后续章节中，我们将说明这个贝尔曼方程如何构成计算、近似和学习 $v_\pi$ 的若干方法的基础。
我们把上面这样的图称为备份图，因为它们描绘了构成强化学习方法核心的更新或备份操作所依据的关系。
这些操作把价值信息从后继状态（或状态--动作对）回传到一个状态（或状态--动作对）。
全书都会使用备份图，为所讨论算法提供图形化概括。
注意，与转移图不同，备份图中的状态节点不一定表示不同状态；例如，一个状态可以是它自己的后继。

\noindent\textbf{例 3.5：网格世界}\quad 图 3.2 左侧展示了一个简单有限 MDP 的矩形网格世界表示。
网格中的格子对应环境的状态。
在每个格子中，有四个可能动作：\texttt{north}、\texttt{south}、\texttt{east} 和 \texttt{west}，它们确定性地使智能体在网格中向相应方向移动一个格子。
会把智能体带出网格的动作，会使其位置保持不变，同时产生 $-1$ 的奖励。
其他动作产生 $0$ 奖励，但从特殊状态 A 和 B 移出时除外。
从状态 A 出发，四个动作都产生 $+10$ 奖励，并把智能体带到 A$'$。
从状态 B 出发，四个动作都产生 $+5$ 奖励，并把智能体带到 B$'$。

\begin{figure}[htbp]
  \centering
  \includegraphics[width=0.66\textwidth]{rl4-gridworld.png}
  \caption{网格世界例子：异常奖励动态（左）与等概率随机策略的状态价值函数（右）。}
\end{figure}

假设智能体在所有状态中都以相等概率选择四个动作。
图 3.2 右侧展示了在 $\gamma=0.9$ 的折扣奖励情形下，该策略的价值函数 $v_\pi$。
这个价值函数是通过求解线性方程组（3.14）得到的。
注意，下边缘附近的值为负；这是因为在随机策略下，那里撞到网格边界的概率较高。
在该策略下，状态 A 是最好的状态，但它的期望回报小于 $10$，也就是小于它的即时奖励，因为从 A 出发智能体会被带到 A$'$，在那里它很可能撞到网格边缘。
另一方面，状态 B 的价值大于 $5$，也就是大于它的即时奖励，因为从 B 出发智能体会被带到 B$'$，而 B$'$ 具有正价值。
从 B$'$ 出发，可能撞到边缘的期望惩罚（负奖励）会被可能偶然到达 A 或 B 的期望收益所抵消，而且抵消有余。

\noindent\textbf{练习 3.14}\quad 贝尔曼方程（3.14）必须对例 3.5 图 3.2 右侧所示的价值函数 $v_\pi$ 的每个状态成立。
请用数值方式说明，对于中心状态（价值为 $+0.7$），相对于它的四个邻近状态（价值分别为 $+2.3$、$+0.4$、$-0.4$ 和 $+0.7$），该方程成立。
这些数值只精确到一位小数。

\noindent\textbf{练习 3.15}\quad 在网格世界例子中，目标的奖励为正，撞到世界边缘的奖励为负，其余时间的奖励为零。
这些奖励的符号重要吗，还是只有它们之间的间隔重要？
请利用（3.8）证明：在所有奖励上加上同一个常数 $c$，会给所有状态的价值都加上同一个常数 $v_c$，因此不会影响任意策略下任意状态之间的相对价值。
$v_c$ 用 $c$ 和 $\gamma$ 如何表示？

\noindent\textbf{练习 3.16}\quad 现在考虑在分幕任务中，例如跑迷宫，对所有奖励都加上同一个常数 $c$。
这会产生任何影响吗，还是会像上面的持续任务那样使任务保持不变？
为什么？
请给出一个例子。

\noindent\textbf{例 3.6：高尔夫}\quad 为了把打一洞高尔夫表述为强化学习任务，我们对每一次击球计入 $-1$ 的惩罚（负奖励），直到把球击入洞中。
状态是球的位置。
一个状态的价值，是从该位置到洞所需击球数的相反数。
我们的动作当然是如何瞄准并击球，以及选择哪支球杆。
我们把前者视为给定，只考虑球杆的选择，并假设球杆要么是推杆，要么是开球杆。
图 3.3 上半部分展示了对总是使用推杆的策略而言，一个可能的状态价值函数 $v_{\text{putt}}(s)$。

\begin{figure}[htbp]
  \centering
  \includegraphics[width=0.58\textwidth]{rl4-golf.png}
  \caption{高尔夫例子：使用推杆时的状态价值函数（上）与使用开球杆时的最优动作价值函数（下）。}
\end{figure}

终止状态“入洞”价值为 $0$。
我们假设从果岭上任何位置都可以推进洞；这些状态的价值为 $-1$。
在果岭外，使用推杆无法到达球洞，因此价值更大。
如果从某个状态通过推杆可以到达果岭，那么该状态的价值必须比果岭的价值少一，也就是 $-2$。
为简单起见，我们假设推杆可以非常精确且确定性地击球，但距离有限。
这就给出了图中标记为 $-2$ 的清晰等值线；这条线与果岭之间的所有位置都正好需要两杆完成该洞。
类似地，位于 $-2$ 等值线的推杆范围内的任何位置都必须有 $-3$ 的价值，如此类推，得到图中所有等值线。
推杆无法使我们脱离沙坑，因此沙坑的价值为 $-\infty$。
总体而言，从发球区用推杆把球打入洞需要六杆。

\noindent\textbf{练习 3.17}\quad 动作价值的贝尔曼方程，即 $q_\pi$ 的贝尔曼方程，是什么？
它必须用状态--动作对 $(s,a)$ 的可能后继的动作价值 $q_\pi(s',a')$ 来表示动作价值 $q_\pi(s,a)$。
提示：右侧备份图对应这个方程。
请给出与（3.14）类似、但针对动作价值的一系列方程。

\begin{figure}[htbp]
  \centering
  \includegraphics[width=0.24\textwidth]{rl4-backup-qpi.png}
  \caption{$q_\pi$ 的备份图。}
\end{figure}

\noindent\textbf{练习 3.18}\quad 一个状态的价值取决于该状态中可能动作的价值，以及在当前策略下每个动作被采取的可能性。
我们可以用一个以该状态为根、并考虑每个可能动作的小型备份图来理解这一点：

\begin{figure}[htbp]
  \centering
  \includegraphics[width=0.62\textwidth]{rl4-ex-3-18-backup.png}
  \caption{练习 3.18 中状态价值与动作价值关系的备份图。}
\end{figure}

请给出与这个直觉和图相对应的方程，用叶节点的期望价值 $q_\pi(s,a)$ 表示根节点的价值 $v_\pi(s)$，其中给定 $S_t=s$。
这个方程应包含一个以遵循策略 $\pi$ 为条件的期望。
然后再给出第二个方程，把期望显式写成关于 $\pi(a|s)$ 的形式，使方程中不再出现期望记号。

\noindent\textbf{练习 3.19}\quad 一个动作的价值 $q_\pi(s,a)$ 取决于期望的下一奖励，以及剩余奖励和的期望。
同样，我们可以用一个小型备份图来理解这一点；这次备份图以一个动作（状态--动作对）为根，并分支到可能的下一状态：

\begin{figure}[htbp]
  \centering
  \includegraphics[width=0.50\textwidth]{rl4-ex-3-19-backup.png}
  \caption{练习 3.19 中动作价值与下一状态价值关系的备份图。}
\end{figure}

请给出与这个直觉和图相对应的方程，用期望下一奖励 $R_{t+1}$ 和期望下一状态价值 $v_\pi(S_{t+1})$ 表示动作价值 $q_\pi(s,a)$，其中给定 $S_t=s$ 和 $A_t=a$。
这个方程应包含一个期望，但不是以遵循策略为条件的期望。
然后给出第二个方程，用（3.2）定义的 $p(s',r|s,a)$ 显式写出期望值，使方程中不再出现期望记号。

\subsection*{3.6\quad 最优策略与最优价值函数}

粗略地说，求解一个强化学习任务，就是找到一个能在长远上获得大量奖励的策略。
对于有限 MDP，我们可以按如下方式精确定义最优策略。
价值函数在策略集合上定义了一个偏序。
如果一个策略 $\pi$ 在所有状态上的期望回报都大于或等于策略 $\pi'$，那么称 $\pi$ 优于或等于 $\pi'$。
换句话说，$\pi\geq\pi'$ 当且仅当对所有 $s\in\mathcal{S}$ 都有 $v_\pi(s)\geq v_{\pi'}(s)$。
总是至少存在一个优于或等于所有其他策略的策略。
这就是最优策略。
虽然最优策略可能不止一个，但我们把所有最优策略都记作 $\pi_*$。
它们共享同一个状态价值函数，称为最优状态价值函数，记作 $v_*$，并定义为
\begin{equation}
v_*(s)\doteq \max_\pi v_\pi(s),
\tag{3.15}
\end{equation}
对所有 $s\in\mathcal{S}$ 成立。

最优策略还共享同一个最优动作价值函数，记作 $q_*$，并定义为
\begin{equation}
q_*(s,a)\doteq \max_\pi q_\pi(s,a),
\tag{3.16}
\end{equation}
对所有 $s\in\mathcal{S}$ 和 $a\in\mathcal{A}(s)$ 成立。
对于状态--动作对 $(s,a)$，这个函数给出的是：在状态 $s$ 中采取动作 $a$，随后遵循一个最优策略时的期望回报。
因此，我们可以用 $v_*$ 表示 $q_*$：
\begin{equation}
q_*(s,a)=\mathbb{E}\left[R_{t+1}+\gamma v_*(S_{t+1})\mid S_t=s,A_t=a\right].
\tag{3.17}
\end{equation}

\noindent\textbf{例 3.7：高尔夫的最优价值函数}\quad 图 3.3 下半部分展示了一个可能的最优动作价值函数 $q_*(s,\texttt{driver})$ 的等值线。
这些值表示：如果我们先用开球杆打一杆，然后在之后每一步选择开球杆或推杆中较好的动作，则每个状态的价值是多少。
开球杆能把球打得更远，但准确性较低。
只有当我们已经非常接近球洞时，才能用开球杆一杆入洞；因此，$q_*(s,\texttt{driver})$ 的 $-1$ 等值线只覆盖果岭的一小部分。
然而，如果我们有两杆，那么就可以从远得多的地方到达球洞，如 $-2$ 等值线所示。
在这种情况下，我们不必把球直接开到小的 $-1$ 等值线以内，只需要开到果岭上的任意位置；从那里我们可以使用推杆。
最优动作价值函数给出的是在承诺采取某个特定的第一个动作之后的价值；在这里这个动作是开球杆，但之后则使用任何最好的动作。
$-3$ 等值线仍然更远，并且包含起始发球区。
从发球区出发，最佳动作序列是两次开球和一次推杆，用三杆把球打入洞中。

因为 $v_*$ 是某个策略的价值函数，所以它必须满足由状态价值贝尔曼方程（3.14）给出的自相容条件。
然而，因为它是最优价值函数，$v_*$ 的相容性条件可以写成一种特殊形式，而不需要引用任何具体策略。
这就是 $v_*$ 的贝尔曼方程，或者称为贝尔曼最优性方程。
直观地说，贝尔曼最优性方程表达的是：某个状态在最优策略下的价值，必须等于从该状态出发采取最佳动作时的期望回报：
\begin{equation}
\begin{aligned}
v_*(s)
&= \max_{a\in\mathcal{A}(s)}q_{\pi_*}(s,a)\\
&= \max_a \mathbb{E}_{\pi_*}[G_t\mid S_t=s,A_t=a]\\
&= \max_a \mathbb{E}_{\pi_*}[R_{t+1}+\gamma G_{t+1}\mid S_t=s,A_t=a]\qquad \text{（由（3.9））}\\
&= \max_a \mathbb{E}[R_{t+1}+\gamma v_*(S_{t+1})\mid S_t=s,A_t=a]\\
&= \max_a \sum_{s',r}p(s',r|s,a)\left[r+\gamma v_*(s')\right].
\end{aligned}
\tag{3.19}
\end{equation}
上面最后两个方程是 $v_*$ 的贝尔曼最优性方程的两种形式。
$q_*$ 的贝尔曼最优性方程为
\begin{equation}
\begin{aligned}
q_*(s,a)
&= \mathbb{E}\left[R_{t+1}+\gamma\max_{a'}q_*(S_{t+1},a')\mid S_t=s,A_t=a\right]\\
&= \sum_{s',r}p(s',r|s,a)\left[r+\gamma\max_{a'}q_*(s',a')\right].
\end{aligned}
\tag{3.20}
\end{equation}

下图中的备份图以图形方式展示了 $v_*$ 和 $q_*$ 的贝尔曼最优性方程所考虑的未来状态和动作范围。
它们与前面给出的 $v_\pi$ 和 $q_\pi$ 的备份图相同，只是在智能体选择点加入了弧线，以表示取的是该选择上的最大值，而不是某个给定策略下的期望值。
左侧备份图以图形方式表示贝尔曼最优性方程（3.19），右侧备份图以图形方式表示（3.20）。

\begin{figure}[htbp]
  \centering
  \includegraphics[width=0.58\textwidth]{rl4-optimal-backups.png}
  \caption{$v_*$ 和 $q_*$ 的备份图。}
\end{figure}

对于有限 MDP，$v_*$ 的贝尔曼最优性方程（3.19）有一个与策略无关的唯一解。
贝尔曼最优性方程实际上是一个方程组，每个状态一个方程；因此如果有 $n$ 个状态，就有 $n$ 个方程和 $n$ 个未知数。
如果环境的动态 $p$ 已知，那么原则上可以使用求解非线性方程组的多种方法之一来求解这个方程组，从而得到 $v_*$。
也可以为 $q_*$ 求解一组相关方程。

一旦有了 $v_*$，确定一个最优策略就相对容易。
对于每个状态 $s$，在贝尔曼最优性方程中使最大值达到的动作会有一个或多个。
任何只把非零概率分配给这些动作的策略都是最优策略。
你可以把这看作一步搜索。
如果你有最优价值函数 $v_*$，那么经过一步搜索后显得最好的动作就是最优动作。
另一种说法是，任何相对于最优评价函数 $v_*$ 贪心的策略都是最优策略。
“贪心”这个术语在计算机科学中用来描述任何只基于局部或即时考虑来选择备选项的搜索或决策过程，而不考虑这种选择可能阻止未来获得更好备选项的可能性。
因此，它描述的是只根据短期后果来选择动作的策略。
$v_*$ 的妙处在于，如果用它来评价动作的短期后果，具体地说是一步后果，那么贪心策略实际上在我们感兴趣的长期意义上是最优的，因为 $v_*$ 已经把所有可能未来行为的奖励后果考虑进去了。
借助 $v_*$，最优的期望长期回报被转化为每个状态上局部且立即可用的量。
因此，一步向前搜索产生的是长期最优动作。

拥有 $q_*$ 会使选择最优动作更加容易。
有了 $q_*$，智能体甚至不必进行一步向前搜索；对任意状态 $s$，它只需找到任意使 $q_*(s,a)$ 最大的动作。
动作价值函数实际上缓存了所有一步向前搜索的结果。
它把最优期望长期回报作为一个对每个状态--动作对局部且立即可用的值给出。
因此，代价是表示一个状态--动作对的函数，而不仅仅是状态的函数；作为回报，最优动作价值函数允许在不知道任何可能后继状态及其价值的情况下选择最优动作，也就是说，不需要知道环境动态的任何信息。

\noindent\textbf{例 3.8：求解网格世界}\quad 假设我们对例 3.5 中引入的简单网格任务求解 $v_*$ 的贝尔曼方程；该任务在图 3.5 左侧再次展示。
回忆一下，状态 A 后跟随 $+10$ 的奖励并转移到状态 A$'$，而状态 B 后跟随 $+5$ 的奖励并转移到状态 B$'$。
图 3.5 中间展示最优价值函数，右侧展示相应的最优策略。
当一个格子中有多条箭头时，所有对应动作都是最优的。

\begin{figure}[htbp]
  \centering
  \includegraphics[width=0.72\textwidth]{rl4-gridworld-optimal.png}
  \caption{网格世界例子的最优解。}
\end{figure}

\noindent\textbf{例 3.9：回收机器人的贝尔曼最优性方程}\quad 使用（3.19），我们可以明确给出回收机器人例子的贝尔曼最优性方程。
为使表达更紧凑，我们分别用 $\mathrm{h}$、$\mathrm{l}$、$\mathrm{s}$、$\mathrm{w}$ 和 $\mathrm{re}$ 缩写状态 \texttt{high} 与 \texttt{low}，以及动作 \texttt{search}、\texttt{wait} 和 \texttt{recharge}。
由于只有两个状态，贝尔曼最优性方程由两个方程组成。
$v_*(\mathrm{h})$ 的方程可以写为
\[
\begin{aligned}
v_*(\mathrm{h})
&=\max\left\{
\begin{array}{l}
p(\mathrm{h}|\mathrm{h},\mathrm{s})[r(\mathrm{h},\mathrm{s},\mathrm{h})+\gamma v_*(\mathrm{h})]
+p(\mathrm{l}|\mathrm{h},\mathrm{s})[r(\mathrm{h},\mathrm{s},\mathrm{l})+\gamma v_*(\mathrm{l})],\\
p(\mathrm{h}|\mathrm{h},\mathrm{w})[r(\mathrm{h},\mathrm{w},\mathrm{h})+\gamma v_*(\mathrm{h})]
+p(\mathrm{l}|\mathrm{h},\mathrm{w})[r(\mathrm{h},\mathrm{w},\mathrm{l})+\gamma v_*(\mathrm{l})]
\end{array}
\right\}\\
&=\max\left\{
\begin{array}{l}
\alpha[r_s+\gamma v_*(\mathrm{h})]+(1-\alpha)[r_s+\gamma v_*(\mathrm{l})],\\
1[r_w+\gamma v_*(\mathrm{h})]+0[r_w+\gamma v_*(\mathrm{l})]
\end{array}
\right\}\\
&=\max\left\{
\begin{array}{l}
r_s+\gamma[\alpha v_*(\mathrm{h})+(1-\alpha)v_*(\mathrm{l})],\\
r_w+\gamma v_*(\mathrm{h})
\end{array}
\right\}.
\end{aligned}
\]
对 $v_*(\mathrm{l})$ 采用同样步骤，得到方程
\[
v_*(\mathrm{l})=\max\left\{
\begin{array}{l}
\beta r_s-3(1-\beta)+\gamma[(1-\beta)v_*(\mathrm{h})+\beta v_*(\mathrm{l})],\\
r_w+\gamma v_*(\mathrm{l}),\\
\gamma v_*(\mathrm{h})
\end{array}
\right\}.
\]
对于任意 $r_s,r_w,\alpha,\beta$ 和 $\gamma$ 的选择，只要 $0\leq\gamma<1$ 且 $0\leq\alpha,\beta\leq1$，都恰好存在一对数 $v_*(\mathrm{h})$ 与 $v_*(\mathrm{l})$ 同时满足这两个非线性方程。

显式求解贝尔曼最优性方程，提供了一条找到最优策略、从而求解强化学习问题的途径。
然而，这个解很少能直接派上用场。
它类似于穷举搜索：向前查看所有可能性，计算它们发生的概率以及它们以期望奖励衡量的可取性。
这个解至少依赖于三个在实践中很少同时成立的假设：（1）我们准确知道环境动态；（2）我们有足够的计算资源来完成该解的计算；（3）马尔可夫性质成立。
对于我们感兴趣的任务类型，一般无法精确实施这个解，因为这些假设会以各种组合被违反。
例如，虽然第一个和第三个假设对双陆棋游戏来说并不构成问题，但第二个假设是主要障碍。
由于该游戏大约有 $10^{20}$ 个状态，即使用当今最快的计算机，求解 $v_*$ 的贝尔曼方程也需要数千年；寻找 $q_*$ 也是如此。
在强化学习中，人们通常必须接受近似解。

许多不同的决策方法都可以看作近似求解贝尔曼最优性方程的方式。
例如，启发式搜索方法可以看作把（3.19）右侧展开若干次，直到某个深度，形成一棵可能性“树”，然后用启发式评价函数近似叶节点上的 $v_*$。
诸如 A$^*$ 的启发式搜索方法几乎总是基于分幕情形。
动态规划方法可以与贝尔曼最优性方程联系得更紧密。
许多强化学习方法则可以清楚地理解为对贝尔曼最优性方程的近似求解，只是用实际经历到的转移来代替对期望转移的知识。
在后续章节中，我们会考察多种这样的方法。

\noindent\textbf{练习 3.20}\quad 画出或描述高尔夫例子的最优状态价值函数。

\noindent\textbf{练习 3.21}\quad 画出或描述高尔夫例子中推杆动作的最优动作价值函数 $q_*(s,\texttt{putter})$ 的等值线。

\noindent\textbf{练习 3.22}\quad 考虑右图所示的持续 MDP。
唯一需要作出的决策位于顶部状态，在那里有两个可用动作：\texttt{left} 和 \texttt{right}。
图中的数字表示每个动作之后确定性地收到的奖励。
恰好有两个确定性策略 $\pi_{\text{left}}$ 和 $\pi_{\text{right}}$。
当 $\gamma=0$ 时，哪个策略最优？
当 $\gamma=0.9$ 时呢？
当 $\gamma=0.5$ 时呢？

\begin{figure}[htbp]
  \centering
  \includegraphics[width=0.32\textwidth]{rl4-ex-3-22-mdp.png}
  \caption{练习 3.22 中的持续 MDP。}
\end{figure}

\noindent\textbf{练习 3.23}\quad 给出回收机器人的 $q_*$ 的贝尔曼方程。

\noindent\textbf{练习 3.24}\quad 图 3.5 给出了网格世界中最佳状态的最优价值，四舍五入到一位小数为 $24.4$。
利用你对最优策略和（3.8）的理解，把这个值符号化地表示出来，然后计算到三位小数。

\noindent\textbf{练习 3.25}\quad 给出一个用 $q_*$ 表示 $v_*$ 的方程。

\noindent\textbf{练习 3.26}\quad 给出一个用 $v_*$ 和四自变量 $p$ 表示 $q_*$ 的方程。

\noindent\textbf{练习 3.27}\quad 给出一个用 $q_*$ 表示 $\pi_*$ 的方程。

\noindent\textbf{练习 3.28}\quad 给出一个用 $v_*$ 和四自变量 $p$ 表示 $\pi_*$ 的方程。

\noindent\textbf{练习 3.29}\quad 用三自变量函数 $p$（3.4）和二自变量函数 $r$（3.5），重写四个价值函数（$v_\pi$、$v_*$、$q_\pi$ 和 $q_*$）的四个贝尔曼方程。

\subsection*{3.7\quad 最优性与近似}

我们已经定义了最优价值函数和最优策略。
显然，学到一个最优策略的智能体已经做得很好，但在实践中这种情况很少发生。
对于我们感兴趣的任务类型，只有付出极大的计算代价才能生成最优策略。
一个定义良好的最优性概念，组织了本书所描述的学习方法，并提供了一种理解各种学习算法理论性质的方式；但它只是一个理想，智能体只能以不同程度去近似它。
如上所述，即使我们拥有环境动态的完整且准确的模型，通常也不可能简单地通过求解贝尔曼最优性方程来计算最优策略。
例如，像国际象棋这样的棋盘游戏只占人类经验的很小一部分，但大型定制计算机仍然无法计算最优走法。
智能体所面对问题的一个关键方面，始终是它可用的计算能力，尤其是它在单个时间步中能够完成的计算量。

可用内存也是一个重要约束。
为了构建价值函数、策略和模型的近似，通常需要大量内存。
在具有小型有限状态集的任务中，可以用数组或表来形成这些近似，每个状态（或状态--动作对）对应一个条目。
我们称这为表格情形，相应方法称为表格方法。
然而，在许多实际感兴趣的情形中，状态数远多于表中可能容纳的条目数。
在这些情况下，必须使用某种更紧凑的参数化函数表示来近似这些函数。

我们对强化学习问题的框定迫使我们接受近似。
不过，它也为获得有用近似提供了一些独特机会。
例如，在近似最优行为时，智能体可能会面对许多概率极低的状态；对这些状态选择次优动作，对智能体接收到的奖励总量几乎没有影响。
例如，Tesauro 的双陆棋程序即使可能会在专家对局中从不出现的棋盘局面上作出很糟糕的决策，仍然能以出色技巧下棋。
事实上，TD-Gammon 可能会在该游戏状态集的很大一部分上作出糟糕决策。
强化学习的在线性质，使得以这样一种方式近似最优策略成为可能：把更多精力放在学习如何为频繁遇到的状态作出好决策上，而以较少关注不常遇到的状态为代价。
这是区分强化学习与其他近似求解 MDP 方法的一个关键性质。

\subsection*{3.8\quad 小结}

现在总结本章中介绍的强化学习问题要素。
强化学习研究的是如何从交互中学习行为，以便达成目标。
强化学习智能体及其环境在一系列离散时间步上交互。
它们接口的规定定义了一个特定任务：动作是智能体作出的选择；状态是作出这些选择的依据；奖励是评价这些选择的依据。
智能体内部的一切都完全为智能体所知且可由它控制；智能体外部的一切都不能被完全控制，但可能完全已知，也可能并非完全已知。
策略是一种随机规则，智能体按它把状态映射为动作选择。
智能体的目标是最大化它随时间接收到的奖励量。

当上述强化学习设置被表述为具有明确定义的转移概率时，它就构成一个马尔可夫决策过程（MDP）。
有限 MDP 是状态、动作和奖励集合（按本章表述）都有限的 MDP。
目前强化学习的大部分理论都限于有限 MDP，但方法和思想适用得更广。

回报是未来奖励的函数，智能体试图最大化其期望值。
它有几种不同定义，取决于任务性质以及是否希望对延迟奖励进行折扣。
无折扣表述适合分幕任务，在这类任务中，智能体--环境交互会自然分解为若干幕；折扣表述适合持续任务，在这类任务中，交互不会自然分解为幕，而是无界持续。
我们试图为这两类任务定义回报，使得一组方程能同时适用于分幕情形和持续情形。

一个策略的价值函数，会把在该策略下从某个状态，或某个状态--动作对，出发的期望回报赋给该状态或状态--动作对。
最优价值函数则把任意策略所能达到的最大期望回报赋给每个状态或状态--动作对。
价值函数为最优的策略就是最优策略。
对于给定 MDP，状态和状态--动作对的最优价值函数是唯一的，但最优策略可能有许多个。
任何相对于最优价值函数贪心的策略都必定是最优策略。
贝尔曼最优性方程是最优价值函数必须满足的特殊相容性条件；原则上，它们可以被求解以得到最优价值函数，并由此相对容易地确定一个最优策略。

强化学习问题可以以许多不同方式提出，取决于关于智能体初始可用知识水平的假设。
在完全知识问题中，智能体拥有环境动态的完整且准确的模型。
如果环境是一个 MDP，那么这样的模型由完整的四自变量动态函数 $p$（3.2）组成。
在不完全知识问题中，环境的完整且完美模型并不可得。

即使智能体拥有完整且准确的环境模型，它通常也无法在每个时间步执行足够多的计算来充分利用它。
可用内存也是一个重要约束。
为了构建价值函数、策略和模型的准确近似，可能需要内存。
在大多数实际感兴趣的情形中，状态数远多于表中可能容纳的条目数，因此必须进行近似。

一个定义良好的最优性概念，组织了本书所描述的学习方法，并提供了一种理解各种学习算法理论性质的方式；但它只是一个理想，强化学习智能体只能以不同程度近似它。
在强化学习中，我们非常关注那些无法找到最优解、而必须以某种方式近似的情形。

\section*{文献与历史说明}

强化学习问题深受最优控制领域中马尔可夫决策过程（MDP）思想的影响。
这些历史影响，以及来自心理学的其他主要影响，在第 1 章给出的简史中已有描述。
强化学习为 MDP 增加了对现实大型问题中的近似和不完全信息的关注。
MDP 与强化学习问题同传统人工智能中的学习和决策问题只存在较弱联系。
不过，人工智能现在正从多种角度积极探索用于规划和决策的 MDP 表述。
MDP 比人工智能先前使用的表述更一般，因为它们允许更一般类型的目标和不确定性。

MDP 理论的处理可见于 Bertsekas（2005）、White（1969）、Whittle（1982，1983）和 Puterman（1994）等。
Ross（1983）给出了有限情形的特别紧凑处理。
MDP 也以随机最优控制为名被研究；其中，自适应最优控制方法与强化学习关系最为密切，例如 Kumar（1985）以及 Kumar 和 Varaiya（1986）。

MDP 理论源自人们理解不确定性下连续决策问题的努力；在这种问题中，每个决策都可以依赖先前决策及其结果。
它有时被称为多阶段决策过程理论，或顺序决策过程理论，并植根于顺序抽样的统计文献；其开端可追溯到 Thompson（1933，1934）和 Robbins（1952）的论文，我们在第 2 章联系赌博机问题时引用过它们。
如果把赌博机问题表述为多情境问题，那么它们就是 MDP 的原型。

据我们所知，最早使用 MDP 形式化讨论强化学习的例子，是 Andreae（1969b）关于学习机器统一观点的描述。
Witten 和 Corbin（1973）实验了一个强化学习系统，后来 Witten（1977，1976a）用 MDP 形式化对其进行了分析。
虽然 Werbos（1977）没有明确提到 MDP，但他提出了与现代强化学习方法相关的随机最优控制问题近似求解方法；另见 Werbos（1982，1987，1988，1989，1992）。
虽然 Werbos 的思想当时没有得到广泛认可，但它们具有先见之明，强调了在包括人工智能在内的多种领域中近似求解最优控制问题的重要性。
强化学习与 MDP 最有影响力的整合归功于 Watkins（1989）。

\noindent\textbf{3.1}\quad 我们用 $p(s',r|s,a)$ 来刻画 MDP 动态的方式稍有不同寻常。
在 MDP 文献中，更常见的是用状态转移概率 $p(s'|s,a)$ 和期望下一奖励 $r(s,a)$ 来描述动态。
然而，在强化学习中，我们更常需要引用单个实际奖励或样本奖励，而不只是它们的期望值。
我们的记号还使得 $S_t$ 和 $R_t$ 通常是联合决定的，因此必须具有相同的时间下标，这一点更加清楚。
在教授强化学习时，我们发现这种记号在概念上更直接，也更容易理解。

关于状态这一系统理论概念的良好直观讨论，可参见 Minsky（1967）。

生物反应器例子基于 Ungar（1990）以及 Miller 和 Williams（1992）的工作。
回收机器人例子受到 Jonathan Connell（1989）建造的收集汽水罐机器人的启发。
Kober 和 Peters（2012）给出了强化学习机器人应用的集合。

\noindent\textbf{3.2}\quad 奖励假设由 Michael Littman 提出（个人交流）。

\noindent\textbf{3.3--4}\quad 分幕任务和持续任务这两个术语不同于 MDP 文献中的通常用法。
在该文献中，通常区分三类任务：（1）有限视界任务，其中交互在某个固定时间步数之后终止；（2）不定视界任务，其中交互可以任意长，但必须最终终止；（3）无限视界任务，其中交互不会终止。
我们的分幕任务和持续任务分别类似于不定视界任务和无限视界任务，但我们更愿意强调交互性质上的差异。
这种差异似乎比通常术语所强调的目标函数差异更根本。
分幕任务常常使用不定视界目标函数，持续任务常常使用无限视界目标函数，但我们把这看作常见的巧合，而不是根本差异。

倒立摆例子来自 Michie 和 Chambers（1968）以及 Barto、Sutton 和 Anderson（1983）。

\noindent\textbf{3.5--6}\quad 基于长期好坏来赋值有古老根源。
在控制理论中，把状态映射到表示控制决策长期后果的数值，是最优控制理论的关键组成部分；该理论在 20 世纪 50 年代通过扩展 19 世纪经典力学中的状态函数理论而发展起来（例如 Schultz 和 Melsa，1967）。
在描述如何给计算机编程以便下棋时，Shannon（1950）建议使用一种评价函数，把棋局位置的长期优势和劣势考虑进去。

Watkins（1989）用于估计 $q_*$ 的 Q-learning 算法（第 6 章）使动作价值函数成为强化学习的重要组成部分，因此这些函数常常被称为“Q 函数”。
不过，动作价值函数的思想比这古老得多。
Shannon（1950）曾建议，国际象棋程序可以使用一个函数 $h(P,M)$ 来决定局面 $P$ 中的走法 $M$ 是否值得探索。
Michie（1961，1963）的 MENACE 系统以及 Michie 和 Chambers（1968）的 BOXES 系统，都可以理解为估计动作价值函数。
在经典物理中，Hamilton 的主函数是一个动作价值函数；牛顿动力学相对于这个函数是贪心的（例如 Goldstein，1957）。
动作价值函数还在 Denardo（1967）用压缩映射处理动态规划的理论中发挥了核心作用。

$v_*$ 的贝尔曼最优性方程由 Richard Bellman（1957a）推广，他称其为“基本函数方程”。
连续时间和连续状态问题中贝尔曼最优性方程的对应物，称为 Hamilton--Jacobi--Bellman 方程（或常简称 Hamilton--Jacobi 方程），这表明其根源在经典物理中（例如 Schultz 和 Melsa，1967）。

高尔夫例子由 Chris Watkins 提出。

\end{document}
